\documentclass[pdftex,14pt,a4paper,russian,utf8]{article}
%\usepackage{pscyr}
\usepackage[utf8]{inputenc}
\usepackage[russian]{babel}
\usepackage[14pt]{extsizes}
\usepackage[pdftex]{graphicx}
\usepackage{geometry}
%\usepackage[T2A]{fontenc}
\usepackage{amssymb,amsmath}
\usepackage{makecell}
\usepackage{eskdpz}
\usepackage{multirow}
\geometry{left=2.5cm}
\geometry{right=1cm}
\geometry{top=2.0cm}
\geometry{bottom=2.5cm}
\linespread{1.3} % 1,5
%\renewcommand{\rmdefault}{ftm}
\setcounter{tocdepth}{2}
\usepackage{graphics}
\usepackage{hyperref}
\usepackage{pdfpages}
\usepackage{wrapfig}
%\renewcommand{\thesection}{}
%\renewcommand{\thesubsection}{}
%\renewcommand{\thesubsubsection}{}
\begin{document}
\begin{titlepage}

\begin{center}
 
\begin{figure}[h!tf]
\begin{center}
\includegraphics[scale=1.5]{./rapidlogo.pdf}\\[1cm]
\end{center}
\end{figure} 

\textsc{\Large Техническое задание}\\[0.5cm]
 
 
% Title
\rule{\linewidth}{0.5mm} \\[0.4cm]
{ \huge { На разработку устройства измерительного для измерения уровней заряда аккумуляторов }}\\[0.4cm]
 
\rule{\linewidth}{0.5mm} \\[1.5cm]

\end{center}

\end{titlepage}



\section{Введение.}
\subsection{Наименование устройства.}
\subsubsection{Наименование устройства : \textit{power board}, далее <<Устройство>>.}
\subsection{Назначение и область применения.}
\subsubsection{Устройство предназначено для работы в составе \textbf{RAPID} с целью.}
\paragraph{}{Измерения уровня заряда трех аккумуляторов и предоставления этой информации основному контролирующему устройству.}
\paragraph{}{Правильной коммутации аккумуляторов в процессе заряда.}
\paragraph{}{Расширения состава датчиков основного контролирующего устройства \textbf{RAPID}}
\section{Требования к устройству.}
\subsection{Требования к функциональным характеристикам.}
\subsubsection{Устройство должно обеспечивать возможность выполнения перечисленных ниже групп функций.}
\paragraph{}{Измерение с помощью аналого-цифровых преобразователей уровней заряда аккумуляторов (напряжений, через делитель напряжения) и опционально (реализуется с помощью перемычек на плате) измерение тока (с помощью токоизмерительного сопротивления и усилителя).}
\paragraph{}{Отключение цепей измерения напряжения и тока и подключения выходов контроллера устройства к выходам на плате напрямую (для сохранения возможности подключения внешних датчиков и исполнительных механизмов).}
\paragraph{}{Коммутация аккумуляторов подключенных к устройству в режиме заряда (о переходе в режим заряда сигнализирует специальный датчик, часть устройства).}
\paragraph{}{Взаимодействие с управляющим устройством посредством линии RS-485 и / или CAN.}
\subsection{Требования к габаритам.}
\subsubsection{Габариты устройства должны вписываться в прямоугольник 60х80 мм.}
\subsection{Требования по унификации.}
\subsubsection{Каждый разъем устройства должен быть уникальным для исключения возможности неправильного подключения.}
\subsection{Требования к надежности.}
\subsubsection{Устройство должно сохранять работоспособность при скачках напряжений аккумуляторов.}
\subsubsection{Устройство должно сохранять работоспособность при питающих напряжениях от 6 до 15 В.}
\end{document}

